\documentclass[11pt]{article}

\usepackage[T1]{fontenc}
\usepackage[margin=1in]{geometry}
\usepackage{amsmath,amssymb,mathtools,amsthm}
\usepackage{booktabs,tabularx,array}
\IfFileExists{lmodern.sty}
  {\usepackage{lmodern}\usepackage{microtype}}
  {\usepackage[expansion=false]{microtype}}
\usepackage{enumitem}
\usepackage[hidelinks]{hyperref}

\setlength{\parindent}{0pt}
\setlength{\parskip}{0.55em}
\setlist[itemize]{leftmargin=*,itemsep=0.15em,topsep=0.25em}
\renewcommand{\arraystretch}{1.12}

\newcommand{\AP}{\operatorname{AP}}
\newcommand{\CNAP}{\operatorname{CNAP}}
\newcommand{\Regime}{\mathcal{R}}
\newcommand{\Assess}{\mathcal{A}}
\newcommand{\CSet}{\mathcal{C}}
\newcommand{\PSet}{\Pi_\star}
\newcommand{\E}{\mathbb{E}}
\newcommand{\Prb}{\mathbb{P}}
\newcommand{\iid}{\mathrel{\overset{\mathrm{iid}}{\sim}}}
\newcommand{\APta}{\widetilde{\AP}}
\newcommand{\CNAPta}{\widetilde{\CNAP}}
\newcommand{\Dta}{\widetilde{\Delta}}
\newcommand{\RobDta}{\underline{\widetilde{\Delta}}_{\PSet}}

\newtheorem{definition}{Definition}
\newtheorem{proposition}{Proposition}
\theoremstyle{remark}
\newtheorem*{remark}{Remark}

\title{Supported Model Replacement Across Target Prevalences\\
\large A Paired Replication-Based Method for Average Precision}
\author{Marcus A. Thomas}
\date{Methodological proposal}

\begin{document}

\maketitle

\begin{abstract}
Average precision depends on the prevalence of the population being evaluated, while a model-replacement claim may fail when the evaluation is repeated. A defensible comparison must therefore preserve pairing while challenging the claimed advantage across both target populations and complete replications. Let $\PSet$ be a declared set of target prevalences, let $\Regime$ be a law for one complete evaluation, and let $K$ be the number of independent evaluations the claim must survive jointly. After prior-standardizing average precision to prevalence $\pi$ and normalizing population random ranking to zero, the proposed replacement estimand is
\[
\theta_{A:B}
=
\E_{\Regime}\left[
\min_{1\leq j\leq K}
\inf_{\pi\in\PSet}
\left\{
\widetilde{\CNAP}_{A,\pi}(D_j)
-
\widetilde{\CNAP}_{B,\pi}(D_j)
\right\}
\right].
\]
The difference is formed on the same units before either challenge is applied, so favorable conditions for one model cannot be matched against unfavorable conditions for the other. The tildes denote average precision averaged over the orderings within each score-tie block. This convention removes expected credit for arbitrary score refinement and, under an exchangeable-label reference law, prevents differing tie structures from creating a positive supported advantage in either orientation. With several independent evaluations, a complete U-statistic estimates the criterion directly. With one evaluation, a paired stratified bootstrap provides only projected support under an empirical fixed-model regime. Replacement requires an outer lower bound for $\theta_{A:B}$ to exceed a prespecified practical margin.
\end{abstract}

\section{What must a replacement claim survive?}
\label{sec:introduction}

Many prediction systems produce a ranking so that a selected set can receive further attention: radiological images are sent for review, candidate molecules are advanced to an assay, and customers are prioritized for contact. In these applications the composition of the selected set matters, not only the separation between outcome classes, because the selected set is what the decision maker receives.

Separation and selected-set purity part company when prevalence changes. Suppose a threshold selects $80\%$ of positive cases and $5\%$ of negative cases. In a balanced population of 10,000 cases, it returns 4,000 positives and 250 negatives, for precision $0.94$. In a population with $1\%$ positives, the same threshold returns 80 positives and 495 negatives, for precision $0.14$. Its true-positive and false-positive rates have not moved. The supply of negative cases has, and that is enough to change the meaning of the selected set.

Average precision aggregates precision as a ranking is traversed, so it inherits the prevalence of the population in which it is evaluated. A study serving one known population can name one target prevalence. A replacement intended for several populations faces a harder choice. Selecting one convenient prevalence hides the populations in which the advantage weakens, whereas averaging requires a defensible distribution over target populations and allows strength in one to compensate for weakness in another. Neither response supports the claim that an advantage holds throughout a declared population set.

Repetition presents the same failure in another direction. An expected advantage can be positive even though one of the evaluations it averages over favors the incumbent model. If replacement requires the advantage to persist under new evaluation units or other declared sources of variation, a stronger result in one evaluation cannot restore a level defeated in another.

The population and replication problems cannot be solved independently. Different evaluations may weaken at different prevalences, so challenging every evaluation at one common prevalence can miss the condition that defeats each of them. A claim required to hold throughout the population set in every evaluation must allow the limiting prevalence to change between evaluations.

Model replacement adds a prior question: which effect should face those challenges? Subtracting two separately challenged model reports allows their adverse cases to occur at different prevalences or in different evaluations. Forming the difference on the same units preserves those conditions, but ordinary finite-sample AP can then reward the model that makes finer score distinctions even when the refinement is unrelated to the outcomes. A replacement effect must preserve pairing without allowing score granularity itself to create an advantage.

The argument follows these failures. It first constructs an AP effect that refers to declared target populations. It then asks what level of a condition-indexed effect survives several complete evaluations without compensation. Only after those challenges have been fixed does the paired comparison return, because the effect entering them must be neutral to arbitrary score refinement. The final sections distinguish direct evidence from independent evaluations from a bootstrap projection based on one evaluation, then state the limits under which either supports replacement.

\section{Comparing performance across target prevalences}
\label{sec:prevalence}

\subsection{Prior-standardized average precision}

Changing the observed positive fraction can alter both the population represented and the evidence used to describe the ranking. Prior standardization separates these operations by changing the class prior while preserving the observed class-conditional score behavior.

At score threshold $c$, let
\[
r(c)=\Prb(S\geq c\mid Y=1),
\qquad
f(c)=\Prb(S\geq c\mid Y=0).
\]
In a target population with prevalence $\pi\in(0,1)$, precision is
\begin{equation}
p_\pi(c)
=
\frac{\pi r(c)}
{\pi r(c)+(1-\pi)f(c)}.
\label{eq:reference-precision}
\end{equation}
Its odds separate the target prior from the ranking behavior:
\begin{equation}
\frac{p_\pi(c)}{1-p_\pi(c)}
=
\frac{\pi}{1-\pi}\frac{r(c)}{f(c)}.
\label{eq:precision-odds}
\end{equation}
Applying Equation~\eqref{eq:reference-precision} across the distinct score thresholds gives $\AP_\pi(D)$, prior-standardized AP for evaluation $D$ at prevalence $\pi$ \cite{elkan2001foundations,siblini2020calibration}. Appendix~\ref{app:ap} gives the empirical definition.

The transformation is licensed only when the class-conditional score distributions transport to the target population. With $S$ denoting a fixed model score, the required condition is
\begin{equation}
P_{\mathrm{obs}}(S\geq c\mid Y=y)
=
P_{\mathrm{target}}(S\geq c\mid Y=y),
\qquad y\in\{0,1\},
\label{eq:score-invariance}
\end{equation}
at the thresholds used by the evaluation. This score-level form of prior-probability shift is weaker than invariance of $P(X\mid Y)$ but follows from it for a fixed scoring rule \cite{saerens2002adjusting,storkey2009training,lipton2018detecting}.

The restriction matters because prevalence can accompany other population changes. A referral population may contain more severe cases within each outcome class, while a new measurement process may change the score for otherwise similar cases. In either setting, Equation~\eqref{eq:reference-precision} reweights rates that need not describe the target. Those changes must be represented by observed evaluations from the corresponding environments rather than by prior adjustment alone \cite{bareinboim2013transportability}.

\subsection{A common effect scale}

Prior-standardized AP still has a different random-ranking value at every target prevalence. Population random ranking has $r(c)=f(c)$ and hence $\AP_\pi=\pi$. Define
\begin{equation}
\CNAP_\pi(D)
=
\frac{\AP_\pi(D)-\pi}{1-\pi}.
\label{eq:cnap}
\end{equation}
CNAP is zero for population random ranking, one for perfect ranking, and negative below chance. The normalization has the chance-corrected form $(\text{observed}-\text{chance})/(\text{ceiling}-\text{chance})$ \cite{cohen1960coefficient}.

Common endpoints do not make the profile prevalence-invariant. At a threshold with true-positive and false-positive rates $r_h$ and $f_h$, the normalized contribution is
\begin{equation}
\frac{p_{\pi,h}-\pi}{1-\pi}
=
\frac{\pi(r_h-f_h)}
{\pi r_h+(1-\pi)f_h}.
\label{eq:threshold-cnap}
\end{equation}
Rare-positive populations therefore charge false positives more heavily. Two models can exchange order as $\pi$ changes even when both profiles improve with prevalence.

Let $\PSet\subset(0,1)$ be a nonempty compact set chosen from deployment plans or scientific knowledge. For a paired profile $\Delta_\pi(D)$, the advantage retained throughout the target set is
\begin{equation}
\underline{\Delta}_{\PSet}(D)
=
\inf_{\pi\in\PSet}\Delta_\pi(D).
\label{eq:robust-difference-generic}
\end{equation}
The infimum answers a noncompensatory question: what is the largest advantage attained at every declared prevalence? A distribution over prevalences would justify an average, but it would support a different claim.

\section{Support across declared conditions and evaluations}
\label{sec:support}

The prevalence profile in Section~\ref{sec:prevalence} is one instance of a condition-indexed effect. More generally, let $\Delta_c(D)$ denote an effect under condition $c\in\CSet$, where every condition in the nonempty set $\CSet$ can be evaluated from one complete evaluation $D$. For AP under score-level prior invariance, $c$ is a target prevalence and $\CSet=\PSet$.

Not every condition can be placed in $\CSet$. If a new site or measurement process changes class-conditional ranking behavior, one evaluation cannot recover its effect by reweighting. Such variation must be observed through complete evaluations generated by $\Regime$. The assessment specification
\[
\Assess=(\CSet,\Regime,K)
\]
therefore separates conditions challenged within each evaluation from mechanisms challenged by generating evaluations; $K$ states how many independent evaluations must retain the effect jointly.

\subsection{The evaluation-generating regime}

A replication is not defined by the act of resampling. It is defined by what the scientific claim permits to vary. Let $\Regime$ denote a probability law for one complete evaluation $D$, and let $D_1,\ldots,D_K\iid\Regime$. Independence is required between complete evaluations, although observations within an evaluation may be clustered or matched when the design and analysis preserve that structure.

The paper's resampling procedure represents the fixed-model regime: the fitted models, environment, and measurement process are held fixed while evaluation units are redrawn. A broader regime can vary environments or repeat development, but it can be estimated only from genuinely corresponding evaluations or from a justified generating hierarchy. Resampling one fixed score vector cannot release mechanisms that the observed evaluation held fixed.

Observed evidence and regime breadth are separate. Several independent evaluations may directly represent either a narrow fixed-model regime or a broader multisite regime. Conversely, a bootstrap may be applied to a study containing several sites, but its scientific reach still depends on which units and levels the resampling scheme redraws.

\subsection{Joint survival}

Expected performance is weaker than performance required to persist. If two independent evaluations yield effects $0.70$ and $0.05$, their mean is $0.375$, although a claim at that level failed in the second evaluation. The greatest level retained by both is their minimum.

\begin{definition}[Replication-survival support]
For $T_1,\ldots,T_K\iid P_{\Regime}^{T}$, where larger values are better, define
\begin{equation}
S_K(T;\Regime)
=
\E_{\Regime}\left[
\min_{1\leq j\leq K}T_j
\right].
\label{eq:support}
\end{equation}
\end{definition}

The realized minimum is the greatest level retained by one set of $K$ evaluations. Its expectation remains on the scale of the effect; it is neither a confidence bound nor the probability that a claim is true. At $K=2$,
\begin{equation}
S_2(T;\Regime)
=
\E_{\Regime}[T]
-
\frac12\E_{\Regime}|T_1-T_2|.
\label{eq:k2}
\end{equation}
The criterion is expected performance minus half the expected disagreement between two independent evaluations.

The interpretation is Popperian only in a limited, negative sense. A level is retained when it survives specified opportunities for defeat, but survival does not establish truth \cite{popper1959logic,popper1963conjectures}. The regime determines the kind of failure presented, while $K$ determines how many independent opportunities are required. Neither choice is supplied by the philosophy.

\subsection{Projected and observed support}
\label{sec:projected}

One observed evaluation has not survived $K$ new evaluations. Resampling it estimates what would survive under an empirical approximation to $\Regime$, so the resulting quantity is \emph{projected support}. The empirical distribution is necessarily closed to score values and failure modes absent from the observed evaluation, which limits the open-ended challenge that motivates the Popperian interpretation.

Genuinely independent evaluations provide a different kind of evidence because they can reveal failures of the regime specification itself. A new site may violate score invariance; a new measurement system may change ranking behavior; a repeated development run may expose instability that a fixed-score bootstrap cannot generate. When such evaluations are available, they enter the estimator directly rather than serving only as an aspiration for future work.

Computational and empirical replication must consequently remain distinct. Increasing the bootstrap count $B$ reduces Monte Carlo error in the same projected quantity. Increasing $K$ changes the criterion by shifting weight toward worse evaluations. Increasing the number of observed independent evaluations changes the evidence available to estimate the criterion.

\subsection{The order of the challenges}

Let
\[
Z_{\CSet}(D)
=
\inf_{c\in\CSet}\Delta_c(D).
\]
The supported advantage is
\begin{equation}
S_{K,\CSet}(\Delta;\Regime)
=
\E_{\Regime}
\left[
\min_{1\leq j\leq K}
\inf_{c\in\CSet}\Delta_c(D_j)
\right].
\label{eq:robust-support}
\end{equation}
At level $s$, the event inside the construction requires $\Delta_c(D_j)>s$ for every $c\in\CSet$ and every $j\leq K$.

Moving the condition search outside support weakens that requirement. The alternative
\[
\inf_{c\in\CSet}S_K(\Delta_c;\Regime)
\]
challenges every evaluation at the same condition before taking the infimum. In the AP application, it can therefore miss a failure at $\pi=0.01$ in one evaluation and a failure at $\pi=0.40$ in another. The intended order satisfies
\begin{equation}
S_K\left(
\inf_{c\in\CSet}\Delta_c;\Regime
\right)
\leq
\inf_{c\in\CSet}
S_K(\Delta_c;\Regime).
\label{eq:order}
\end{equation}
Equality holds when the same condition minimizes the relevant profile almost surely; Appendix~\ref{app:support-proofs} gives the proof.

The challenge becomes stronger when either index set grows.

\begin{proposition}[Challenge monotonicity]
\label{prop:monotonicity}
If $\mathcal C_1\subseteq\mathcal C_2$ and $K_1\leq K_2$, then
\[
S_{K_2,\mathcal C_2}(\Delta;\Regime)
\leq
S_{K_1,\mathcal C_1}(\Delta;\Regime).
\]
\end{proposition}

No comparable ordering holds for arbitrary changes of $\Regime$, because releasing environments or development runs changes the law of the effect rather than enlarging an index set within the same law.

\section{Why replacement must be paired and tie-neutral}
\label{sec:paired}

\subsection{Marginal summaries do not compose}

A replacement decision concerns an advantage realized under the same conditions. Let
\[
A_{\PSet}(D)
=
\inf_{\pi\in\PSet}\CNAP_{A,\pi}(D),
\qquad
B_{\PSet}(D)
=
\inf_{\pi\in\PSet}\CNAP_{B,\pi}(D).
\]
Subtracting their supported marginal values allows the two models to meet different adverse prevalences and different adverse evaluations.

The error is one-sided:
\begin{equation}
S_K(A_{\PSet};\Regime)
-
S_K(B_{\PSet};\Regime)
\geq
S_K\left(
\inf_{\pi\in\PSet}
\{\CNAP_{A,\pi}-\CNAP_{B,\pi}\};
\Regime
\right).
\label{eq:noncomposition}
\end{equation}
The proof in Appendix~\ref{app:paired-proofs} uses only properties of the infimum and minimum. Separate reports therefore overstate the supported paired advantage by an amount they do not identify.

\subsection{Ordinary pairing rewards arbitrary refinement}

Pairing aligns the units and outcomes, but ordinary finite-sample AP still depends on the score vectors' tie structures. Stepwise AP is upward displaced under random ordering because it evaluates precision when positive cases enter the ranking \cite{bestgen2015random}. A fully tied score vector has one threshold and returns the population random-ranking value, while a vector of distinct scores creates more opportunities for favorable positive entries.

Four units show why this matters for replacement. Let model $A$ assign four distinct scores and model $B$ assign one common score. With two positive labels, two negative labels, and $\pi=1/2$, assign the labels uniformly over the six possibilities. Neither score vector is associated with the outcomes. Under the ordinary threshold-block convention, model $A$ has expected AP $49/72$, while model $B$ has AP $1/2$ under every assignment. Hence
\[
\E[\CNAP_A-\CNAP_B]
=
\frac{13}{36}.
\]
For two independent assignments, the supported difference remains positive at $29/216$. Pairing has favored the model that broke ties without information.

A null anchor cannot be subtracted uniformly because the displacement depends on model performance. It is large near no association but vanishes at perfect separation, where both models reach the ceiling. Appendix~\ref{app:paired-proofs} gives an exact four-unit example in which centering by the joint-permutation anchor favors the coarser model even though both models separate the classes perfectly. The defect must therefore be removed before the nonlinear population and replication challenges are applied.

\subsection{Tie-averaged average precision}
\label{sec:tie-average}

Let $\mathcal O(D)$ be the strict rankings consistent with a score vector, obtained by ordering tied units in every possible way while preserving the order of the score blocks.

\begin{definition}[Tie-averaged average precision]
\begin{equation}
\APta_\pi(D)
=
\frac{1}{|\mathcal O(D)|}
\sum_{o\in\mathcal O(D)}
\AP_\pi(D_o),
\label{eq:tie-average}
\end{equation}
where $D_o$ is evaluation $D$ resolved into strict ranking $o$. Define $\CNAPta_\pi$ through Equation~\eqref{eq:cnap}.
\end{definition}

The convention is equivalent to adding independent continuous jitter that reorders units only within their original tie blocks, then averaging ordinary AP over that outcome-blind resolution. Because the model supplied no ordering within a block, the comparison does not credit whichever arbitrary ordering happened to align with the labels. Appendix~\ref{app:tie-computation} evaluates the average in closed form without enumeration or random jitter.

Tie averaging still rewards informative refinement. A finer score improves realized AP when it orders cases correctly within a coarse block and loses AP when it orders them incorrectly. Only refinement unrelated to the outcomes receives zero expected credit.

\begin{proposition}[Neutrality to arbitrary refinement]
\label{prop:refinement}
Suppose model $A$ preserves the order of model $B$'s score blocks and, conditionally on $S_B$ and $Y$, its strict ranking is uniform over the rankings consistent with $S_B$. Then, for every $\pi$,
\[
\E\!\left[
\APta_{A,\pi}\mid S_B,Y
\right]
=
\APta_{B,\pi}.
\]
\end{proposition}

The proposition conditions on the realized outcomes, so it applies when the coarse model is informative. It removes expected credit for arbitrary refinement without requiring exchangeable evaluation units.

\subsection{The supported paired advantage}

The paired profile is
\[
\Dta_\pi(D)
=
\CNAPta_{A,\pi}(D)
-
\CNAPta_{B,\pi}(D),
\]
and the advantage retained within evaluation $D$ is
\[
\RobDta(D)
=
\inf_{\pi\in\PSet}\Dta_\pi(D).
\]
The model-replacement estimand is
\begin{equation}
\theta_{A:B}
=
S_K(\RobDta;\Regime)
=
\E_{\Regime}
\left[
\min_{1\leq j\leq K}
\inf_{\pi\in\PSet}
\{\CNAPta_{A,\pi}(D_j)-\CNAPta_{B,\pi}(D_j)\}
\right].
\label{eq:paired-estimand}
\end{equation}
The difference is formed before either minimum. A positive value is therefore an advantage for $A$ retained throughout the declared prevalence set under the joint-replication criterion.

\subsection{Reference behavior and design scope}
\label{sec:paired-scope}

Refinement neutrality compares a fine ranking with a coarser ranking it refines. Two arbitrary uninformative rankings require a stronger reference law.

\begin{definition}[Exchangeable-label reference law]
\label{def:exchangeable}
Conditionally on both score vectors and the positive count $n_+$, every label vector containing $n_+$ positives has probability $\binom{n}{n_+}^{-1}$.
\end{definition}

Under this law, labels read along any fixed strict ranking are uniform over the same set of arrangements. Expected tie-averaged AP therefore depends on $n_+$, $n_-$, and $\pi$, but not on the score vector or its tie structure.

\begin{proposition}[Reference behavior of paired support]
\label{prop:paired-reference}
Under the exchangeable-label reference law,
\[
\E_0[\Dta_\pi]=0
\qquad\text{for every }\pi\in\PSet,
\]
and consequently
\[
S_K(\RobDta;\Regime_0)\leq0
\]
in both orientations.
\end{proposition}

The inequality need not be equality because the prevalence and replication minima can move an uninformative paired effect below zero. The result prevents differing tie structures from manufacturing a positive supported advantage, although the challenge itself may be conservative.

Full exchangeability is stronger than independence. Matched, clustered, and heterogeneous-risk designs usually permit only a restricted group of label arrangements, under which score vectors aligned with the strata can have different reference values. Tie averaging still removes arbitrary ordering within score blocks, and Equation~\eqref{eq:paired-estimand} remains descriptive, but the no-positive-anchor result no longer follows. Appendix~\ref{app:paired-proofs} gives a matched-pair counterexample.

\subsection{No verdict and replacement}

The reversed orientation is
\[
\theta_{B:A}
=
S_K\left(
\inf_{\pi\in\PSet}\{-\Dta_\pi\};
\Regime
\right).
\]
Because the support operator contains minima, the reversed value is not generally the negative of the original. The two orientations nevertheless satisfy
\begin{equation}
\theta_{A:B}+\theta_{B:A}\leq0.
\label{eq:orientation}
\end{equation}
Both models therefore cannot receive a positive supported advantage over the other.

The inequality permits a no-verdict region. When both orientations are nonpositive, the declared challenges defeat superiority in either direction. At $K=2$, the width of
\[
\left[\theta_{A:B},-\theta_{B:A}\right]
\]
is
\[
\E\!\left[
\max_{j\leq2}\sup_{\pi\in\PSet}\Dta_\pi(D_j)
-
\min_{j\leq2}\inf_{\pi\in\PSet}\Dta_\pi(D_j)
\right].
\]
It therefore contains both variation across target prevalences and disagreement across evaluations. When $\PSet$ is a singleton, the width reduces to the expected absolute disagreement between the paired effects of two evaluations.

Let $d\geq0$ be the smallest retained CNAP advantage that would justify replacing model $B$ by model $A$. The performance criterion is
\begin{equation}
\theta_{A:B}>d.
\label{eq:replacement}
\end{equation}
A point estimate does not establish this inequality. The proposed decision rule requires a validated outer lower bound for $\theta_{A:B}$ to exceed $d$, with pairing preserved through every estimation layer.

\section{Estimation}
\label{sec:estimation}

\subsection{Several observed evaluations}

Suppose $R$ genuinely independent complete evaluations $D_1,\ldots,D_R\iid\Regime$ are available. For each evaluation, compute
\[
z_r
=
\inf_{\pi\in\PSet}
\{\CNAPta_{A,\pi}(D_r)-\CNAPta_{B,\pi}(D_r)\}.
\]
When $R\geq K$, the complete order-$K$ U-statistic \cite{hoeffding1948ustat}
\begin{equation}
\widehat\theta_{A:B}^{\,\mathrm{obs}}
=
\binom RK^{-1}
\sum_{\substack{I\subseteq\{1,\ldots,R\}\\|I|=K}}
\min_{r\in I}z_r
\label{eq:observed-u}
\end{equation}
is unbiased for Equation~\eqref{eq:paired-estimand}. Sorting avoids enumerating all subsets; Appendix~\ref{app:computation} gives the resulting formula.

Observed evaluations make the inferential basis transparent, but they do not remove the need to define $\Regime$. Several test sets from one measurement system identify a narrower regime than test sets spanning sites or repeated development runs. When $R$ is small, the U-statistic can also be imprecise, especially as $K$ approaches $R$ and few distinct subsets remain.

\subsection{One observed evaluation}

When only one independent evaluation is available, a paired stratified bootstrap can project the fixed-model criterion. Let $D_{\mathrm{obs},+}$ and $D_{\mathrm{obs},-}$ contain the positive and negative evaluation units. For computational replication $b=1,\ldots,B$, sample the declared numbers of units with replacement within each class, keep both model scores attached to every sampled unit, and compute
\[
z_b
=
\inf_{\pi\in\PSet}
\{\CNAPta_{A,\pi}(D^{(b)})-\CNAPta_{B,\pi}(D^{(b)})\}.
\]
Applying the complete order-$K$ U-statistic to $z_1,\ldots,z_B$ estimates support under the empirical within-class score distributions.

Stratification conditions the projected regime on its positive and negative counts. A same-design analysis uses the observed counts, whereas a prospective analysis uses the counts planned for future evaluations. If the claim instead allows class counts to vary, their generating law belongs in $\Regime$. Clustered or matched studies must resample the coarsest unit treated as independent, which may preclude simple class-stratified sampling.

The same sampled units must be used for both models and throughout the prevalence profile. Resampling models separately would destroy pairing, while resampling separately at each prevalence would prevent one evaluation from revealing its own limiting population.

The bootstrap represents only variation available in the observed score pools. It cannot generate unseen score values or failures caused by a new environment, measurement process, or development run. Its claim must therefore be reported as projected fixed-model support rather than as survival of $K$ observed replications.

\subsection{Tail fidelity, prevalence search, and sparse outcomes}

The role of $K$ creates an estimation tension. From Appendix~\ref{app:support-proofs},
\[
S_K(T;\Regime)
=
\int_0^1Q_T(v)K(1-v)^{K-1}\,dv,
\]
so increasing $K$ moves weight toward the lower tail of the evaluation-effect distribution. The stronger challenge consequently depends more heavily on whether the observed evaluations, or their empirical bootstrap distribution, represent adverse outcomes. This dependence is especially consequential when positive outcomes are sparse and the distribution of AP is coarse.

A finite set $\PSet=\{\pi_1,\ldots,\pi_G\}$ defines an exact finite challenge. A grid used to approximate an interval can miss an interior minimum and raise the reported value. A one-dimensional optimizer can instead search a compact interval, with its tolerance checked against a dense grid.

Nonunique or moving prevalence minimizers require a further distinction. Uniformly accurate profiles may still yield an accurate infimum value even when the minimizing prevalence is unstable. Inference about that value is less regular, however, because the infimum is nonsmooth at flat or multiple minima. Consistency of the projected point estimate, fidelity of its lower tail, and coverage of an outer bound must therefore be studied separately.

\subsection{Outer uncertainty}

An outer resampling layer estimates uncertainty from the observed evidence. With several evaluations, the independent evaluations form the outer resampling units. With one evaluation, units are resampled according to the study design before the inner fixed-model projection is repeated. Every outer sample must repeat pairing, tie averaging, the prevalence search, and the order-$K$ support estimator.

Ordinary bootstrap coverage is not automatic because AP can be discrete and the minimizing prevalence can change under perturbation. Percentile, basic, studentized, and bias-corrected bounds may behave differently for the same design. The replacement rule therefore treats the lower bound as a procedure requiring design-specific simulation or theory rather than assigning nominal coverage by convention.

\section{Interpretation and reporting}
\label{sec:discussion}

\subsection{What the number supports}

The estimand is an assessment-relative magnitude, not a portable property of either model. It reports the expected paired advantage retained throughout $\PSet$ by all $K$ evaluations generated under $\Regime$. Omitting any of those declarations changes what the number means.

Projected and observed estimates require different language. An estimate based on several genuinely independent evaluations may be described as replicated evidence under the declared regime. A one-evaluation bootstrap estimates projected support under empirical within-class score distributions and should not be described as having survived external replications.

The result also remains a ranking-performance claim. It does not evaluate calibration, costs, fairness, feasibility, or the consequences of acting on the selected cases. Those considerations can determine deployment even when Equation~\eqref{eq:replacement} is satisfied.

\subsection{Minimum reporting content}

A concise report should identify the claim before presenting diagnostics. Table~\ref{tab:reporting} gives the minimum content.

\begin{table}[ht]
\centering
\caption{Minimum reporting content for a supported replacement analysis.}
\label{tab:reporting}
\begin{tabularx}{\textwidth}{@{}p{0.23\textwidth}X@{}}
\toprule
Component & Report \\
\midrule
Supported quantity &
$\widehat\theta_{A:B}$, the reversed orientation, the margin $d$, and the outer lower bound used for the decision \\
Assessment &
$\PSet$ and its scientific basis; $\Regime$; what is fixed and redrawn; the independent evaluation unit; and $K$ \\
Evidence &
Whether complete evaluations are observed or bootstrap-projected; observed class counts; replication counts; and resampling hierarchy \\
Population transport &
The score-level invariance used to evaluate target prevalences and the populations to which it is asserted to apply \\
Comparison &
Model orientation, paired sampling, tie-averaged convention, and the basis for the exchangeable-label reference law \\
Diagnostics &
Paired-difference profile, distribution of minimizing prevalences, evaluation-effect distribution, and sensitivity to $K$ and $\PSet$ \\
Computation &
Prevalence grid or optimizer, bootstrap count where applicable, Monte Carlo error, and interval method \\
\bottomrule
\end{tabularx}
\end{table}

A possible headline statement is:
\begin{quote}
Across the declared prevalence set $\PSet$ and $K$ independent evaluations under the fixed-model regime, model $A$ had an estimated tie-averaged supported CNAP advantage of [value] over model $B$; the reversed orientation was [value]. The outer lower bound [did or did not] exceed the prespecified replacement margin $d=[value].$
\end{quote}
When the estimate comes from one test set, ``independent evaluations'' must be replaced by ``projected independent evaluations under an empirical fixed-model regime.''

\subsection{Limits and metric reach}

Prior standardization does not transport ranking behavior across changes within outcome class. If severity, measurement, or feature distributions change $P(S\mid Y)$, target-population evaluations are required.

Tie averaging does not create full label exchangeability. In matched, clustered, or heterogeneous-risk designs, the paired estimand remains computable but may inherit a design-specific anchor unless the admissible comparisons respect the design's symmetry.

AUROC separates the general challenge structure from the complications specific to AP. Pure prior shift leaves AUROC unchanged, so a prevalence set supplies no additional challenge while the class-conditional score distributions remain fixed. An AUROC claim can nevertheless fail across conditions that change ranking behavior within class. Averaging across those conditions reintroduces the same compensation problem, and the support operator applies when the conditions can be evaluated within each complete evaluation; otherwise they must be represented through $\Regime$.

The apparent simplification does not make the present estimator portable to AUROC without further work. Half-credit for score ties gives AUROC an exchangeable-label expectation of $1/2$ regardless of tie structure, which removes the positive reference anchor found for AP. The supported paired difference still depends on its null dispersion, however, so different score structures can produce unequal negative values in the two orientations. Nor can the relevant AUROC conditions usually be generated by prior reweighting: subgroup minima divide the available sample, while environmental conditions require corresponding data. A full AUROC application must therefore solve a different estimation problem.

Finally, the choices of $\PSet$, $\Regime$, and $K$ remain scientific commitments. Widening $\PSet$ or increasing $K$ mechanically lowers the criterion, but neither operation proves that the resulting assessment captures the failures relevant to deployment.

\section{Conclusion}

A model-replacement claim based on average precision must survive more than a favorable marginal summary. The population must be named because AP changes with prevalence, and the evaluations must be named because an expected advantage can conceal a failure to persist.

Prior standardization evaluates each model at the target prevalences when class-conditional score behavior is invariant. Chance normalization places those values on a common population random-to-perfect scale. Tie averaging then removes expected credit for arbitrary refinement before the models are compared on the same units. The prevalence infimum and replication minimum are applied only after that paired difference has been formed:
\[
\theta_{A:B}
=
\E_{\Regime}
\left[
\min_{j\leq K}
\inf_{\pi\in\PSet}
\{\CNAPta_{A,\pi}(D_j)-\CNAPta_{B,\pi}(D_j)\}
\right].
\]

Several independent evaluations estimate this criterion directly through a U-statistic. One evaluation can support only a bootstrap projection under an empirical fixed-model regime. That projection can guide a replacement analysis, but it cannot substitute computational resampling for new empirical opportunities to reveal failure.

\appendix

\section{Empirical prior-standardized average precision}
\label{app:ap}

\subsection{Definition and weighted implementation}

Let
\[
D=\{(s_i,y_i)\}_{i=1}^{n},
\qquad
y_i\in\{0,1\},
\]
with $n_+=\sum_i y_i>0$ and $n_-=n-n_+>0$. Let $c_1>\cdots>c_H$ be the distinct observed score values. At threshold $h$, define
\[
\operatorname{TP}_h
=
\sum_i\mathbf 1\{y_i=1,s_i\geq c_h\},
\qquad
\operatorname{FP}_h
=
\sum_i\mathbf 1\{y_i=0,s_i\geq c_h\},
\]
and
\[
r_h=\frac{\operatorname{TP}_h}{n_+},
\qquad
f_h=\frac{\operatorname{FP}_h}{n_-}.
\]
With $r_0=0$, noninterpolated prior-standardized AP under the threshold-block convention is
\begin{equation}
\AP_\pi(D)
=
\sum_{h=1}^{H}
(r_h-r_{h-1})
\frac{\pi r_h}
{\pi r_h+(1-\pi)f_h}.
\label{eq:empirical-ap}
\end{equation}
All units sharing a score enter one threshold block. A block containing only negative cases has zero recall increment, although its false positives reduce the precision of later blocks.

The same quantity can be computed by weighting each positive unit by $w_+=\pi/n_+$ and each negative unit by $w_-=(1-\pi)/n_-$. Weighted precision at threshold $h$ is Equation~\eqref{eq:reference-precision}, while weighted recall remains $r_h$.

\subsection{Behavior across prevalence}

Combining Equations~\eqref{eq:cnap} and \eqref{eq:empirical-ap} gives
\begin{equation}
\CNAP_\pi(D)
=
\sum_{h=1}^{H}
(r_h-r_{h-1})
\frac{\pi(r_h-f_h)}
{\pi r_h+(1-\pi)f_h}.
\label{eq:cnap-profile}
\end{equation}
The derivative of a threshold contribution is
\begin{equation}
\frac{\partial}{\partial\pi}
\left\{
\frac{\pi(r_h-f_h)}
{\pi r_h+(1-\pi)f_h}
\right\}
=
\frac{f_h(r_h-f_h)}
{\{\pi r_h+(1-\pi)f_h\}^2}.
\label{eq:cnap-derivative}
\end{equation}
If $r_h\geq f_h$ at every threshold with positive recall increment, the profile is nondecreasing and its minimum over $[\pi_L,\pi_U]$ occurs at $\pi_L$. A paired difference can remain nonmonotone even when both marginal profiles are nondecreasing.

\section{Properties of replication-survival support}
\label{app:support-proofs}

\subsection{Survival and quantile representations}

For a random variable $T\in[L,U]$,
\[
T=L+\int_L^U\mathbf 1\{T>s\}\,ds.
\]
For $K$ independent copies, expectation of their minimum gives
\begin{equation}
S_K(T;\Regime)
=
L+\int_L^U\Prb_{\Regime}(T>s)^K\,ds.
\label{eq:survival-representation}
\end{equation}
At each level $s$, the powered survival probability is the probability that all $K$ evaluations retain the level.

If $Q_T$ is the quantile function of $T$, the smallest of $K$ independent uniform variables has density $K(1-v)^{K-1}$. Hence
\begin{equation}
S_K(T;\Regime)
=
\int_0^1Q_T(v)K(1-v)^{K-1}\,dv.
\label{eq:quantile-weight}
\end{equation}
The weight moves toward the lower tail as $K$ increases.

\subsection{Proof of the operator-order inequality}

Let $X_c^{(j)}=\Delta_c(D_j)$ and $Y_c=\min_{j\leq K}X_c^{(j)}$. For every realized array,
\[
\min_j\inf_{c\in\CSet}X_c^{(j)}
=
\inf_{c\in\CSet}\min_jX_c^{(j)}
=
\inf_{c\in\CSet}Y_c.
\]
Since $\inf_c Y_c\leq Y_{c'}$ for every $c'\in\CSet$,
\[
\E\left[\inf_{c\in\CSet}Y_c\right]
\leq
\inf_{c\in\CSet}\E[Y_c],
\]
which proves Equation~\eqref{eq:order}.

\subsection{Proof of challenge monotonicity}

Take $\mathcal C_1\subseteq\mathcal C_2$ and $K_1\leq K_2$. Draw $K_2$ evaluations once and use the first $K_1$ for the smaller assessment. For every realized array,
\[
\min_{j\leq K_2}\inf_{c\in\mathcal C_2}\Delta_c(D_j)
\leq
\min_{j\leq K_1}\inf_{c\in\mathcal C_1}\Delta_c(D_j).
\]
Expectation preserves the inequality and proves Proposition~\ref{prop:monotonicity}.

\section{Paired comparison proofs and counterexamples}
\label{app:paired-proofs}

\subsection{Why marginal subtraction overstates paired advantage}

Let $U_\pi$ and $V_\pi$ be two profiles, with
\[
U_{\PSet}=\inf_{\pi\in\PSet}U_\pi,
\qquad
V_{\PSet}=\inf_{\pi\in\PSet}V_\pi,
\qquad
T_{\PSet}=\inf_{\pi\in\PSet}(U_\pi-V_\pi).
\]
Choosing a prevalence minimizing $U_\pi$ gives
\[
T_{\PSet}
\leq
U_{\PSet}-V_\pi
\leq
U_{\PSet}-V_{\PSet}.
\]
Monotonicity of $S_K$ therefore gives
\[
S_K(U_{\PSet}-V_{\PSet})
\geq
S_K(T_{\PSet}).
\]

For effects evaluated on the same replications, the minimum is superadditive:
\[
\min_j(X_j+Y_j)
\geq
\min_jX_j+\min_jY_j.
\]
Taking $X=U_{\PSet}-V_{\PSet}$ and $Y=V_{\PSet}$ yields
\[
S_K(U_{\PSet})-S_K(V_{\PSet})
\geq
S_K(U_{\PSet}-V_{\PSet}),
\]
which completes Equation~\eqref{eq:noncomposition}.

\subsection{Proof of refinement neutrality}

Condition on $S_B$ and $Y$. Proposition~\ref{prop:refinement} makes model $A$'s strict ranking $O_A$ uniform on $\mathcal O(D_B)$. Tie-averaged AP for $A$ is the conditional expectation of ordinary AP over whatever ties $A$ retains. By the tower property,
\[
\E[\APta_{A,\pi}\mid S_B,Y]
=
\E[\AP_\pi(D_{O_A})\mid S_B,Y].
\]
Uniformity of $O_A$ turns the right side into the average of $\AP_\pi$ over $\mathcal O(D_B)$, which is $\APta_{B,\pi}$.

\subsection{Rank invariance under exchangeable labels}

Fix $n_+$, $n_-$, and a strict ranking $o$. Under Definition~\ref{def:exchangeable}, the label sequence read along $o$ is uniform over the $\binom n{n_+}$ arrangements containing $n_+$ positives. Reading along another strict ranking is a bijection of the same arrangements, so the sequence has the same law.

For a strict ranking, $\pi\mapsto\CNAP_\pi(D_o)$ is a deterministic function of that sequence. Its expectation is therefore the same for every strict ranking. Tie-averaged AP is a convex combination of strict-ranking values, so any two score vectors on the same units satisfy
\[
\E_0[\CNAPta_{A,\pi}]
=
\E_0[\CNAPta_{B,\pi}].
\]
Thus $\E_0[\Dta_\pi]=0$.

Finally,
\[
\E_0[\RobDta]
\leq
\inf_{\pi\in\PSet}\E_0[\Dta_\pi]
=0,
\]
and $S_K(T)\leq\E[T]$. Their composition proves Proposition~\ref{prop:paired-reference}; replacing $\Dta$ by $-\Dta$ proves the reversed orientation.

\subsection{Why restricted label symmetry is insufficient}

Take four units in two matched pairs, each containing exactly one positive. The design permits labels to swap within pairs but not between pairs. Let model $A$ give all units the same score, while model $B$ places the first pair above the second and ties the units within each pair. At $\pi=1/2$, exact tie averaging over the four permitted assignments gives
\[
\CNAPta_A-\CNAPta_B
=
\frac1{36}
\]
in one orientation for every assignment. The effect has no replication variation, so its supported value remains $1/36$ for every $K$.

The obstruction is the symmetry group of the design. Full exchangeability places every strict ranking in one orbit of the label-permutation group. Restricted designs may divide rankings into several orbits with different anchors, and averaging within score ties cannot move a ranking between them.

\subsection{Why paired null centering fails}

Take two positive and two negative units at $\pi=1/2$, with
\[
s_A=(4,3,2,1),
\qquad
s_B=(2,2,1,1),
\]
and let the two highest-scored units be positive. Both models separate the classes perfectly, so every stratified bootstrap evaluation has CNAP one for each model and paired difference zero.

Under joint label permutation, exact enumeration of the nested $K=2$ procedure gives a design anchor of $35/384$ for $A-B$ and $-295/1152$ for $B-A$. Subtracting these anchors would report a positive corrected value of about $0.256$ for the second orientation, favoring the coarser model although both predictions are perfect. The displacement is not an additive constant because perfect separation leaves no room for it.

\subsection{Proof of the orientation bound}

For evaluation $j$, write
\[
L_j=\inf_{\pi\in\PSet}\Dta_\pi(D_j),
\qquad
U_j=\sup_{\pi\in\PSet}\Dta_\pi(D_j).
\]
The two orientations satisfy
\[
\theta_{A:B}
=
\E\left[\min_{j\leq K}L_j\right],
\qquad
\theta_{B:A}
=
-\E\left[\max_{j\leq K}U_j\right].
\]
Since $L_j\leq U_j$ for every evaluation,
\[
\theta_{A:B}+\theta_{B:A}
=
\E\left[
\min_{j\leq K}L_j-\max_{j\leq K}U_j
\right]
\leq0,
\]
which proves Equation~\eqref{eq:orientation}. The negative of this sum is the width of the no-verdict region. When $\PSet$ is a singleton and $K=2$, $L_j=U_j=T_j$, so the width reduces to $\E|T_1-T_2|$.

\section{Closed-form computation of tie-averaged AP}
\label{app:tie-computation}

Index score blocks $j=1,\ldots,J$ from highest to lowest. Let block $j$ contain $a_j$ positive and $b_j$ negative units, with $m_j=a_j+b_j$, and define
\[
A_j=\sum_{\ell<j}a_\ell,
\qquad
B_j=\sum_{\ell<j}b_\ell.
\]
Then
\begin{equation}
\APta_\pi(D)
=
\sum_{j=1}^{J}
\frac{a_j}{n_+}\frac1{m_j}
\sum_{k=1}^{m_j}
\sum_{t=\max(0,k-1-b_j)}^{\min(k-1,a_j-1)}
\frac{\binom{a_j-1}{t}\binom{b_j}{k-1-t}}
{\binom{m_j-1}{k-1}}
\frac{\pi r_{jkt}}
{\pi r_{jkt}+(1-\pi)f_{jkt}},
\label{eq:tie-closed}
\end{equation}
where
\[
r_{jkt}=\frac{A_j+t+1}{n_+},
\qquad
f_{jkt}=\frac{B_j+k-1-t}{n_-}.
\]

To obtain the formula, fix one positive unit in block $j$. Under a uniformly random ordering of the block, its position $k$ is uniform on $\{1,\ldots,m_j\}$. Given $k$, the number $t$ of other positive units in the $k-1$ preceding positions is hypergeometric with $a_j-1$ available positives and $b_j$ negatives. The final fraction in Equation~\eqref{eq:tie-closed} is the prior-standardized precision when that positive enters the ranking. Averaging over $t$ and $k$, multiplying by the block's $a_j$ positive units, and summing over blocks gives the result.

The calculation costs $O(\sum_jm_j^2)$ operations. Its combinatorial weights do not depend on $\pi$ and can be reused throughout the prevalence search. When all scores are distinct, Equation~\eqref{eq:tie-closed} reduces to ordinary stepwise AP.

\section{Computation and outer uncertainty}
\label{app:computation}

\subsection{General support order}

Sort $R$ observed evaluation effects:
\[
z_{(1)}\leq\cdots\leq z_{(R)}.
\]
The value $z_{(i)}$ is the minimum in every order-$K$ subset containing it and $K-1$ of the $R-i$ larger values. Hence Equation~\eqref{eq:observed-u} equals
\begin{equation}
\widehat\theta_{A:B}^{\,\mathrm{obs}}
=
\binom RK^{-1}
\sum_{i=1}^{R-K+1}
\binom{R-i}{K-1}z_{(i)}.
\label{eq:sorted-observed}
\end{equation}

For $B$ bootstrap effects, the projected estimator is
\begin{equation}
\widehat\theta_{A:B}^{\,\mathrm{proj}}
=
\binom BK^{-1}
\sum_{i=1}^{B-K+1}
\binom{B-i}{K-1}z_{(i)}.
\label{eq:sorted-bootstrap}
\end{equation}
At $K=2$, this reduces to
\[
\widehat\theta_{A:B}^{\,\mathrm{proj}}
=
\frac{2}{B(B-1)}
\sum_{i=1}^{B-1}(B-i)z_{(i)}.
\]

\subsection{Monte Carlo error}

For $K=2$, define the leave-one-replication projection
\[
\widehat h_{1,-i}(z_i)
=
\frac1{B-1}\sum_{j\ne i}\min(z_i,z_j).
\]
A first-order Monte Carlo standard error for the projected estimator is
\begin{equation}
\widehat{\operatorname{se}}_{\mathrm{MC}}
=
\frac2{\sqrt B}
\operatorname{sd}_{i}
\left\{
\widehat h_{1,-i}(z_i)
\right\}.
\label{eq:mcse}
\end{equation}
Prefix sums over the sorted array compute all projections in linear time after sorting.

\subsection{Prevalence evaluation and outer bounds}

For a finite prevalence set with $G$ values, score-block counts are computed once in each evaluation and reused across prevalences. For a continuous interval, Equation~\eqref{eq:cnap-profile} is a smooth one-dimensional objective away from thresholds with no recall contribution. Endpoint values and stationary points may be checked directly, while a numerical optimizer should be validated against a dense grid.

With several observed evaluations, an outer sample resamples the complete evaluations. With one observed evaluation, an outer sample resamples the independent units or clusters according to the study design, and the inner bootstrap then projects complete fixed-model evaluations from that outer empirical distribution. Pairing must be retained in both layers, and every outer sample must repeat the prevalence search. Because the target combines a lower-tail functional with a potentially nonsmooth infimum, coverage should be established under designs resembling the intended application.

\begin{thebibliography}{99}

\bibitem{bareinboim2013transportability}
Bareinboim, E., \& Pearl, J. (2013).
A general algorithm for deciding transportability of experimental results.
\emph{Journal of Causal Inference}, 1(1), 107--134.

\bibitem{bestgen2015random}
Bestgen, Y. (2015).
Exact expected average precision of the random baseline for system evaluation.
\emph{The Prague Bulletin of Mathematical Linguistics}, 103, 131--138.

\bibitem{cohen1960coefficient}
Cohen, J. (1960).
A coefficient of agreement for nominal scales.
\emph{Educational and Psychological Measurement}, 20, 37--46.

\bibitem{elkan2001foundations}
Elkan, C. (2001).
The foundations of cost-sensitive learning.
In \emph{Proceedings of the Seventeenth International Joint Conference on Artificial Intelligence}, 973--978.

\bibitem{hoeffding1948ustat}
Hoeffding, W. (1948).
A class of statistics with asymptotically normal distribution.
\emph{Annals of Mathematical Statistics}, 19, 293--325.

\bibitem{lipton2018detecting}
Lipton, Z. C., Wang, Y.-X., \& Smola, A. (2018).
Detecting and correcting for label shift with black box predictors.
In \emph{Proceedings of the 35th International Conference on Machine Learning}, PMLR 80, 3122--3130.

\bibitem{popper1959logic}
Popper, K. R. (1959).
\emph{The Logic of Scientific Discovery}.
Hutchinson.

\bibitem{popper1963conjectures}
Popper, K. R. (1963).
\emph{Conjectures and Refutations: The Growth of Scientific Knowledge}.
Routledge.

\bibitem{saerens2002adjusting}
Saerens, M., Latinne, M., \& Decaestecker, C. (2002).
Adjusting the outputs of a classifier to new a priori probabilities.
\emph{Neural Computation}, 14(1), 21--41.

\bibitem{siblini2020calibration}
Siblini, W., Fr\'ery, J., He-Guelton, L., Obl\'e, F., \& Wang, Y. (2020).
Master your metrics with calibration.
In \emph{Advances in Intelligent Data Analysis XVIII}, 457--469.

\bibitem{storkey2009training}
Storkey, A. (2009).
When training and test sets are different: characterizing learning transfer.
In \emph{Dataset Shift in Machine Learning}, 3--28. MIT Press.

\end{thebibliography}

\end{document}
